\section{Related Work}
\label{sec:related}

\para{Loopholes in humans and machines.} The loophole framing originates in cognitive science studies of children's compliance behavior~\citep{bridgers2025loophole, qian2024loophole}: a loophole requires recovering the speaker's intent and then choosing a different but linguistically valid interpretation. \citet{murthy2023loopholes} introduced the first LLM evaluation, finding that only GPT-3.5+ class models reproduce the human pattern across the trouble, upset, and humor axes, and that no LLM differentiates humor---suggesting the higher-order theory-of-mind reasoning behind funny loopholes is absent. \citet{choi2025loopholes} provide the first direct measurement of \emph{exploitation} (rather than judgment) in modern LLMs across scalar implicature, bracketing ambiguity, and the 108-stimulus power-relationship corpus from \citet{bridgers2025loophole}. Their three main findings ground our setup: (i) intent prediction is near-perfect on multiple choice ($\sim$100\%), so exploitation is strategic, not parsing failure; (ii) exploitation rates scale with capability for non-reasoning models; (iii) long-chain-of-thought distillations (DeepSeek-R1 family) do not exploit loopholes the way matched non-reasoning models do. We adopt their paired-inverse template trick and their hyperparameters (T$=0.7$, top-$p=0.95$, 3 seeds) so our numbers sit on the same scale, but extend the regime from goal conflict to knowledge conflict.

\para{Honesty, refusal, and knowledge boundaries.} \citet{li2024honesty} define honesty as recognizing what one knows and faithfully expressing it; \citet{cheng2024dontknow} and \citet{chen2024knboundary} build supervised signals to teach LLMs to refuse on out-of-knowledge queries. \citet{yin2023selfaware} release \selfaware---answerable/unanswerable question pairs spanning subjective, future, and counterfactual cases---which we repurpose here as a behavioral probe rather than a refusal-training target. \citet{cao2023learnrefuse} treats knowledge-scope-limitation plus refusal as the primary mitigation lever. Our \kcl benchmark complements all of the above: existing work measures refusal versus answer; \kcl adds the missing middle option (\dodge) and the missing causal control (paired knowable items).

\para{Hallucination as a decision, not a parameter readout.} \citet{simhi2025choke} introduce CHOKE---a hallucination class in which the model can answer correctly under one phrasing but confidently fabricates under another---arguing that hallucination is partly a decision-stage failure, not a pure knowledge failure. Internal-state probes~\citep{azaria2023internal, orgad2024llmsknow, kossen2024semantic, duan2024hidden} converge on the same picture: hidden states encode ``this output is wrong'' even when the surface output is confident. This is mechanistically central for us. If knowledge-conflict loopholes were a knowledge-only phenomenon, the model would simply not have the relevant tokens and would either refuse or produce uncorrelated noise. The CHOKE and probe evidence implies a deliberate output-channel exploit, which our \emph{confident hallucination} finding in \kcl is consistent with. \citet{zhang2023sirens} survey the broader hallucination space; \citet{zhao2024curriculum} name the laziness-versus-hallucination trade-off in which our ``obviously unsatisfactory'' middle zone lives.

\para{Safety-training loopholes.} \citet{wei2023jailbroken} frame jailbreak success as either \emph{competing objectives} or \emph{mismatched generalization}---both loophole patterns. \citet{andriushchenko2024pasttense} show that simply rephrasing a harmful request in the past tense bypasses refusal training in most frontier models, a striking near-trivial loophole. \citet{denison2024subterfuge} demonstrate zero-shot generalization across a sycophancy-to-reward-tampering curriculum, suggesting simple and sophisticated loopholes live on a continuum. We treat these as evidence that loophole behavior is not hypothetical for frontier models, but they all operate in the goal-conflict regime our work complements.

\para{Pragmatic instruction following.} \citet{zhixuan2024clips} propose CLIPS, a Bayesian agent that jointly infers user goals and instruction noise---the cooperative inverse of a loophole exploiter. \citet{liu2023ambiguity} and \citet{stengel2024zero} study whether LLMs detect ambiguity at all; their answer (often no) is one reason loophole exploitation can succeed---if ambiguity is unrecognized, the model picks an interpretation without flagging the choice.

\para{Chain-of-thought faithfulness.} \citet{turpin2023cot} and \citet{barez2025cot} show that verbalized reasoning does not always reflect the model's actual computation. We inherit this caveat for our \rseventy results: when the reasoning trace is the visible response, surface honesty signals are downstream of an inspection that may itself be unfaithful.

\para{Adjacent risks.} \citet{park2024deception} survey deceptive behaviors more broadly; \citet{greenblatt2024faking} document alignment faking under training-time observation; \citet{sharma2023sycophancy} document sycophancy as a base-rate specification-gaming behavior. \citet{laine2024sad} caution that frontier models can detect when they are being evaluated, which may attenuate dodge rates relative to deployment. None of these directly measure knowledge-conflict loopholes, but each names a related way in which surface compliance can mask intent violation.

\para{What we add.} \kcl is the first benchmark that (i) pairs each unknowable question with a topic-matched knowable control so loophole rate is measurable causally and (ii) operationalizes a four-category response taxonomy whose middle terms (\dodge, \hedged) are precisely what existing refusal-versus-answer scoring drops. Combined with the \selfaware repurposing and the \choi single-model replication, this lets us test whether the dodge mechanism reported by \citet{choi2025loopholes} for goal conflict appears in the knowledge-conflict regime---and, as we will show, the answer is: not as the dominant mode, but with a clear hallucination-shaped analog.
